% =====================================================
% 题型：三角恒等变换与三角形最值解答题 (q_lyb_01_18_trig_proof_max.tex)
% =====================================================
% =====================================================
% 难度：★★★ (难题)
% =====================================================
\newcommand{\qTriangleTrigProofMaxStem}{(1) 在单位圆中，用向量方法证明两角差的余弦公式：\\
\[ \cos(\alpha - \beta) = \cos\alpha\cos\beta + \sin\alpha\sin\beta; \]\\
(2) 若 \(\cos(\alpha - \beta) = \dfrac{3}{5}\)，\(\cos(\alpha + \beta) = \dfrac{1}{5}\)，求 \(\tan\alpha\tan\beta\) 的值;\\
(3) 在 \(\triangle ABC\) 中，求 \(\cos A + \cos B + \cos C\) 的最大值。}

\newcommand{\qTriangleTrigProofMaxAnswer}{
(1) 见解析; (2) \(\dfrac{1}{2}\); (3) \(\dfrac{3}{2}\).
}

\newcommand{\qTriangleTrigProofMaxSolution}{%
  (1) 证明：
  在平面直角坐标系中，以原点 \(O\) 为圆心画单位圆。
  设角 \(\alpha\) 与 \(\beta\) 的终边分别与单位圆交于点 \(A\) 和 \(B\)。
  根据三角函数的几何定义，点 \(A\) 和 \(B\) 的坐标分别为：
  \[ A(\cos\alpha, \sin\alpha), \quad B(\cos\beta, \sin\beta) \]
  则向量 \(\vec{OA} = (\cos\alpha, \sin\alpha)\)，\(\vec{OB} = (\cos\beta, \sin\beta)\)。
  一方面，根据向量数量积的坐标表示：
  \[ \vec{OA} \cdot \vec{OB} = \cos\alpha\cos\beta + \sin\alpha\sin\beta \]
  另一方面，根据数量积的定义有：
  \[ \vec{OA} \cdot \vec{OB} = |\vec{OA}| \cdot |\vec{OB}| \cos\langle\vec{OA}, \vec{OB}\rangle \]
  因为 \(A, B\) 在单位圆上，所以 \(|\vec{OA}| = |\vec{OB}| = 1\)。
  两向量 \(\vec{OA}\) 与 \(\vec{OB}\) 夹角的大小可表示为 \(|\alpha - \beta|\)（或 \(\alpha - \beta\) 的某种周期表示）。
  由于余弦函数是偶函数且周期为 \(2\pi\)，则：
  \[ \cos\langle\vec{OA}, \vec{OB}\rangle = \cos(\alpha - \beta) \]
  由此可得：
  \[ \vec{OA} \cdot \vec{OB} = 1 \cdot 1 \cdot \cos(\alpha - \beta) = \cos(\alpha - \beta) \]
  因此：
  \[ \cos(\alpha - \beta) = \cos\alpha\cos\beta + \sin\alpha\sin\beta \]
  两角差的余弦公式得证。\\
  
  (2) 解：
  因为：
  \[ \cos(\alpha - \beta) = \cos\alpha\cos\beta + \sin\alpha\sin\beta = \frac{3}{5} \quad \text{--- (i)} \]
  \[ \cos(\alpha + \beta) = \cos\alpha\cos\beta - \sin\alpha\sin\beta = \frac{1}{5} \quad \text{--- (ii)} \]
  将 (i) 与 (ii) 两式相加，得：
  \[ 2\cos\alpha\cos\beta = \frac{4}{5} \implies \cos\alpha\cos\beta = \frac{2}{5} \]
  将 (i) 与 (ii) 两式相减，得：
  \[ 2\sin\alpha\sin\beta = \frac{2}{5} \implies \sin\alpha\sin\beta = \frac{1}{5} \]
  因此：
  \[ \tan\alpha\tan\beta = \frac{\sin\alpha\sin\beta}{\cos\alpha\cos\beta} = \frac{\frac{1}{5}}{\frac{2}{5}} = \frac{1}{2} \]
  即 \(\tan\alpha\tan\beta\) 的值为 \(\dfrac{1}{2}\)。\\
  
  (3) 解：
  在 \(\triangle ABC\) 中，有 \(A + B + C = \pi\)，即 \(C = \pi - (A+B)\)。
  我们有：
  \[ \cos A + \cos B + \cos C = 2\cos\frac{A+B}{2}\cos\frac{A-B}{2} + \cos C \]
  因为 \(\cos C = 1 - 2\sin^2\frac{C}{2}\)，且 \(\cos\frac{A+B}{2} = \cos\frac{\pi-C}{2} = \sin\frac{C}{2}\)，
  所以：
  \[ \cos A + \cos B + \cos C = 2\sin\frac{C}{2}\cos\frac{A-B}{2} + 1 - 2\sin^2\frac{C}{2} \]
  整理为关于 \(\sin\frac{C}{2}\) 的二次函数式：
  \[ \cos A + \cos B + \cos C = -2\sin^2\frac{C}{2} + 2\cos\frac{A-B}{2}\sin\frac{C}{2} + 1 \]
  因为 \(A-B \in (-\pi, \pi)\)，所以 \(\cos\frac{A-B}{2} \le 1\)。
  又因 \(\sin\frac{C}{2} > 0\)，可得：
  \[ \cos A + \cos B + \cos C \le -2\sin^2\frac{C}{2} + 2\sin\frac{C}{2} + 1 \]
  配方得：
  \[ -2\sin^2\frac{C}{2} + 2\sin\frac{C}{2} + 1 = -2\left(\sin\frac{C}{2} - \frac{1}{2}\right)^2 + \frac{3}{2} \]
  当且仅当 \(\sin\frac{C}{2} = \frac{1}{2}\)（即 \(C = \dfrac{\pi}{3}\)）且 \(\cos\frac{A-B}{2} = 1\)（即 \(A = B = \dfrac{\pi}{3}\)）时，等号成立。
  此时三角形为等边三角形。
  所以 \(\cos A + \cos B + \cos C\) 的最大值为 \(\dfrac{3}{2}\)。%
}

\newcommand{\qTriangleTrigProofMaxDiagnosis}{%
  学生第 (2) 问正确。第 (1) 问证明两角差余弦公式时不够严谨，存在用相关结论反推的风险；第 (3) 问知道最大值为 \(\frac32\)，但证明偏口头化，缺少完整不等式推导。%
}

\newcommand{\qTriangleTrigProofMaxTeachingNotes}{%
  精讲。第 (1) 问按单位圆坐标与数量积证明；第 (3) 问补半角配方或凹性证明，避免只记“等边最大”的结论。%
}


% =====================================================
% 别名兼容：旧课次宏定义
% =====================================================
\newcommand{\qLYBZeroOneEighteenStem}{\qTriangleTrigProofMaxStem}
\newcommand{\qLYBZeroOneEighteenAnswer}{\qTriangleTrigProofMaxAnswer}
\newcommand{\qLYBZeroOneEighteenSolution}{\qTriangleTrigProofMaxSolution}
\newcommand{\qLYBZeroOneEighteenDiagnosis}{\qTriangleTrigProofMaxDiagnosis}
\newcommand{\qLYBZeroOneEighteenTeachingNotes}{\qTriangleTrigProofMaxTeachingNotes}
